\section{Related work}
\label{sec:related}

\paragraph{Egocentric human data for robot learning.}
Egocentric video has long served as a representation prior: Ego4D~\cite{grauman2022ego4d} supplies large-scale first-person video, and R3M/VIP-style pretraining~\cite{nair2022r3m,ma2022vip} distills it into visual features for downstream control.
A more recent line injects egocentric data \emph{directly into policy training}: EgoMimic~\cite{kareer2024egomimic} uses hand-mounted grippers whose kinematics match the robot, UMI~\cite{chi2024universal} uses hand-held grippers with visual-inertial pose recovery, DexCap~\cite{wang2024dexcap} uses portable mocap, and EgoBridge~\cite{egobridge2025} studies mixing in simulation with explicit domain gaps.
Cross-embodiment robot corpora~\cite{openx2024} face a cousin of the same problem across robot morphologies.
These works demonstrate practical benefits, substitution ratios~\cite{zevaego2025}, and failure modes, and motivate our complementary regression-level analysis of how task mismatch and sample scale affect an explicitly optimized mixing weight.

\paragraph{Domain adaptation and transfer regression.}
Classical domain-adaptation theory bounds target error by source error plus a discrepancy and a shared-optimal-hypothesis term~\cite{bendavid2010theory}; subsequent work studies when quantity compensates for quality, the value of target data, and when representation invariance can hurt~\cite{bendavid2014quantity,hanneke2019value,zhao2019learning}.
More closely, \cite{yang2025precise} analyze two-task linear regression under covariate and model shifts in the proportional high-dimensional limit. They characterize the bias and variance of hard parameter sharing, identify positive-to-negative transfer transitions, and study source rebalancing.
\cite{kovacevic2026private} consider weighted joint ridge regression with multiple heterogeneous sources, deriving deterministic risk equivalents that capture covariance shift, model shift, sample sizes, regularization, and privacy noise. Their formulation contains pooled weighted ridge as a limiting case and includes Bayes-optimal tuning under Gaussian priors.
These works already provide precise transfer-risk analyses and source-weighting principles under task mismatch. We focus on the global geometry along a scalar mixing weight: a minimum-dimension multivalley construction, certification of all population candidates, and finite-sample selection based on a uniform fixed-dimension risk expansion.

\paragraph{Mixture scaling laws and the present scope.}
\cite{wu2026mixture} study mixed ridge regression with a shared regression function, heterogeneous covariances and noise levels, and derive minimax and ridge scaling laws under aligned power-law spectra. Sim-to-real co-training analyses provide complementary empirical models of multi-source learning~\cite{lei2026simreal}.
Relative to this line, we allow an explicit deterministic task mismatch $\delta$ and study the geometry and selection of optimal mixing weights at fixed dimension. Scaling exponents in a proportional regime are outside the present analysis.

\paragraph{Technical tools.}
Our analysis uses ridge-regression risk decompositions~\cite{dobriban2018high,hastie2022surprises}, with target risk measured in the $\Sigma_r$ geometry. Real Wishart moments and resolvent expansions retain the joint sample-Gram fluctuations of full ridge. The per-coordinate label-feedthrough calculation belongs to the separate diagonalized estimator in the supplementary. Invariant-representation limitations~\cite{zhao2019learning,arjovsky2019invariant} provide related context; alignment and gradient-conflict experiments are not part of the current main-text validation.
